\chapter{General Conclusion and Proposed Method for the PFE}

\section{General Conclusion}
This Master thesis studied Digital Twin technology in the context of eHealth. It showed that Digital Twins can support the transition from fragmented digital health applications toward dynamic, personalized, and service-oriented healthcare systems. A Digital Twin is not only a digital copy or a visualization tool; it is a connected virtual representation that uses data, models, analytics, simulation, and services to support monitoring, interpretation, prediction, and decision-making.

The study first clarified the conceptual foundations of Digital Twins, including the distinction between digital models, digital shadows, Digital Twins, and intelligent Digital Twins. It then introduced the eHealth and digital-health environment in which these concepts must be adapted. The healthcare-oriented review showed that Digital Twins for Health are receiving growing attention in personalized medicine, remote monitoring, hospital processes, public health, and cardiovascular applications. However, the field remains emerging, and many works are still conceptual, prototype-based, or limited to specific scenarios.

The comparative analysis produced a synthesis of the main findings from the structured literature review. It derived a general multi-layer eHealth Digital Twin architecture, identified cross-cutting requirements, described the data and feedback flow, and proposed evaluation criteria for trustworthy systems. The logic of the thesis is therefore: general Digital Twin foundations, eHealth context, healthcare applications, structured literature comparison, synthesized architecture and evaluation requirements, cardiovascular deep-dive, and finally the CardioTwin PFE specialization.

\section{Main Contributions of the Master Study}
The main contributions of this Master study are summarized as follows:

\begin{itemize}[leftmargin=*]
    \item clarification of Digital Twin concepts, maturity levels, and enabling technologies;
    \item analysis of eHealth foundations, including connected devices, health data, interoperability, and digital-health services;
    \item review of healthcare Digital Twin applications, with attention to cardiovascular examples;
    \item comparative study of selected general, healthcare-oriented, architecture-oriented, ethical, and cardiovascular works;
    \item synthesis of a general layered architecture for eHealth Digital Twins from recurring elements in the literature;
    \item identification of evaluation requirements covering technical, model, clinical, workflow, ethical, and governance dimensions;
    \item identification of research gaps related to synchronization, interoperability, validation, explainability, security, privacy, and clinical integration;
    \item preparation of a theoretical foundation for the CardioTwin PFE.
\end{itemize}

\section{Limitations of the Study}
This thesis is mainly a literature-based state-of-the-art and conceptual synthesis. It does not provide a full clinical deployment, a certified medical device, or a prospective clinical validation. The terminology of the field is heterogeneous, and published works do not always use the term Digital Twin with the same level of maturity. Some systems described in the literature are closer to digital models, simulations, or digital shadows than to fully synchronized and clinically integrated Digital Twins.

Another limitation is the limited number of mature real-world healthcare Digital Twin deployments. The field evolves rapidly, and the synthesis depends on the evidence available in the reviewed literature. Finally, the general architecture derived in Chapter 6 still requires use-case-specific implementation, data selection, model design, regulatory analysis, and validation before it can support clinical practice.

\section{From the Master Study to the PFE}
The Master thesis studied Digital Twins for eHealth from a general perspective. It analyzed healthcare applications and architectures, extracted general architectural and evaluation findings, and identified the requirements for trustworthy Digital Twin-based systems. The PFE builds on this foundation by operationalizing part of the synthesis in one clinical domain. It should therefore be understood as a practical specialization of the Master findings, not as proof that all unresolved research gaps have been solved.

The selected specialization is cardiovascular health. The cardiovascular domain satisfies the technical and methodological criteria identified in the literature synthesis: it requires longitudinal and personalized monitoring; it provides electrical, hemodynamic, clinical, physiological, and imaging information; patient presentations are heterogeneous; disease trajectories change over time; and treatment response varies between patients~\cite{coorey2022health,corralacero2020digital}. The reviewed literature provides examples ranging from numerical simulation to patient-specific calibrated models and AI-assisted interpretation~\cite{geddes2025right,gu2025heartfailure}. This makes the domain a realistic pathway for an academic prototype while still requiring careful clinical validation before real medical deployment.

\section{Introduction to the PFE Project}
The proposed PFE is titled \textit{CardioTwin -- Design and Deployment of an AI-Powered Cardiovascular Digital Twin for Diagnostic Assistance and Therapeutic Simulation}. CardioTwin is introduced as a practical cardiovascular specialization of the general eHealth Digital Twin architecture synthesized in this Master study.

CardioTwin focuses on a subset of the general architecture derived in Chapter 6. It does not attempt to implement every layer and requirement at full clinical scale. Instead, it concentrates on a patient-oriented cardiovascular prototype in which ECG signals, systolic blood pressure, diastolic blood pressure, heart rate, demographic information, clinical history, symptoms, and other physiological measurements can provide complementary views of the patient's state. Clinicians may benefit from integrated diagnostic assistance, longitudinal follow-up, risk summaries, and controlled what-if simulations that explore possible therapeutic scenarios.

The objective of the PFE is to design and implement a patient-oriented cardiovascular Digital Twin platform that combines real biomedical measurements, ECG signal analysis, clinical and physiological data, AI-based diagnostic assistance, patient-state representation, cardiovascular risk assessment, therapeutic what-if simulation, monitoring, and visualization services.

\section{Proposed PFE Methodology}
The proposed PFE methodology follows a sequence derived from the general architecture and adapted to the cardiovascular use case.

\subsection{Real-World Data Acquisition}
The first step consists of acquiring cardiovascular and clinical data. The data may include ECG signals, systolic blood pressure, diastolic blood pressure, heart rate when available, patient demographic information, clinical history, symptoms, and other relevant physiological measurements. In the practical prototype, connected medical sensors and the MySignals eHealth development kit may be used to support biomedical data acquisition, depending on availability and implementation constraints.

\subsection{Data Preprocessing and Validation}
Acquired data must be prepared before they can update the twin or feed diagnostic modules. This step includes ECG filtering, artifact handling, normalization, segmentation, measurement validation, missing-data handling, and timestamp alignment. The purpose is to reduce noise, identify unreliable measurements, and preserve the temporal coherence required for longitudinal monitoring.

\subsection{ECG Diagnostic-Assistance Module}
The ECG module uses deep-learning models to analyze ECG signals for rhythm or abnormality detection. Biomedical datasets such as PTB-XL, MIT-BIH, SVDB, AFDB, or INCART may be used during development or evaluation, depending on the selected tasks and available resources. The goal is to provide diagnostic assistance, not autonomous diagnosis.

\subsection{Clinical and Physiological Assessment}
The system also integrates broader patient information, including age, sex, systolic blood pressure, diastolic blood pressure, clinical history, symptoms, and relevant cardiac measurements. These variables support a more complete cardiovascular assessment than ECG classification alone.

\subsection{Patient-State Fusion}
CardioTwin should fuse ECG findings, clinical risk indicators, physiological measurements, and historical information into a unified virtual representation of the patient's cardiovascular state. This makes the twin more than an isolated ECG classifier. The patient state should preserve the distinction between measured values, calculated values, model-estimated values, predicted values, and simulated values.

\subsection{Digital Twin Update}
The cardiovascular twin is updated when new validated patient information is received. Measured values may come directly from sensors or clinical records. Calculated values may be produced from formulas or preprocessing steps. Model-estimated values may be inferred from incomplete or indirect observations. Predicted values may represent expected future trends. Simulated values may represent hypothetical scenarios. These categories must not be confused because they have different levels of certainty and clinical meaning.

\subsection{What-If Therapeutic Simulation}
The what-if module is an exploratory clinical decision-support component. It may simulate changes in blood pressure, modifications of physiological parameters, lifestyle scenarios, treatment or medication-effect scenarios, and possible evolution over a selected time horizon. The generated scenarios are decision-support estimates based on models, evidence, or population-level effects and must not be interpreted as guaranteed individual clinical outcomes or autonomous medical prescriptions.

\subsection{Service and Visualization Layer}
The platform should provide patient dashboards, ECG visualization, historical trends, risk summaries, alerts, what-if scenario outputs, and clinician-oriented reports. These services must present information clearly and must distinguish observations, predictions, and simulations.

\subsection{Human-in-the-Loop Decision Support}
CardioTwin assists healthcare professionals but does not replace them. Clinical decisions remain under professional supervision. Simulated recommendations require medical interpretation, and autonomous therapeutic control is outside the scope of the PFE.

\section{Evaluation Strategy}
The PFE evaluation strategy should cover the main components of the prototype. Data acquisition evaluation will verify correct sensor communication, completeness, timestamps, signal quality, and measurement consistency. ECG-model evaluation will use metrics such as accuracy, precision, recall, F1-score, macro-F1, sensitivity, specificity, confusion matrices, and external dataset testing when available.

Patient-state evaluation will examine consistency between measurements and the stored state, correct longitudinal updates, and the handling of missing or contradictory measurements. What-if evaluation will focus on traceability of assumptions, consistency with cited medical evidence, plausible direction and magnitude of simulated effects, transparent uncertainty, warnings, and the absence of unsupported claims about individual treatment efficacy.

System evaluation will consider latency, modularity, reliability, usability, security, interoperability, and dashboard clarity. Full clinical evaluation would require clinician review, prospective assessment, ethical approval where necessary, representative patient data, and regulatory consideration. Such full validation remains outside the guaranteed scope of the academic prototype and should be treated as future work unless it is explicitly performed.

\section{Future Work}
Future work should first focus on implementing the CardioTwin prototype and integrating real biomedical sensors. Further work should improve the fusion of ECG findings with clinical and physiological data, validate patient-state updates, strengthen what-if simulations, add explainability and uncertainty communication, and integrate interoperability standards such as HL7 FHIR where appropriate.

Additional research should include clinician-centered usability evaluation, stronger cybersecurity and consent mechanisms, prospective validation, and a clearer path toward regulatory and ethical assessment. In a broader perspective, the same methodology could later be extended to other healthcare Digital Twin domains beyond cardiovascular monitoring.

\section{Final Remarks}
Digital Twins for eHealth represent a promising research direction at the intersection of healthcare, computer science, data engineering, artificial intelligence, simulation, and systems modeling. Their success will depend not only on technical performance, but also on trust, interoperability, clinical relevance, explainability, and responsible design. CardioTwin provides a practical continuation of this Master study by applying the general findings to a cardiovascular Digital Twin prototype.

\clearpage
